\section{软件虚拟化：QEMU的全系统实现}\label{sec:software}

本章考察纯软件方式构造整台机器的路径：以 \code{qemu-system-riscv64} 的 TCG 全系统模拟为对象，guest 为 RISC-V，host 为 x86-64，机器采用 \code{virt} 平台的 CLINT 与 PLIC 配置，不涉及 RISC-V H 扩展或 AIA\cite{source3}。这一配置下，宿主处理器不执行任何 guest 指令，guest 的 CPU、内存、中断控制器与外围设备全部由运行在宿主用户态的一个 QEMU 进程以软件数据结构和方法调用来实现。QEMU 用翻译代码和软件例程实现 guest 的状态转移。CPU 指令翻译只是其中一环，地址空间、中断和设备必须共同提供一致的整机接口。

\subsection{CPU虚拟化：RISC-V到TCG IR再到x86-64}\label{sec:software-cpu}

先看全景。QEMU 的 CPU 执行由译时与运行时两条路径组成，代码块~\ref{lst:tcg-loop} 给出伪代码。运行时主循环以 PC 与翻译相关 CPU 状态查 TB 缓存，命中直接执行已缓存的宿主代码，未命中才现场翻译；译时路径把 guest 指令逐条解码、翻译成 TCG IR，译到控制转移封块，再经优化与寄存器分配发射宿主机器码。

\begin{lstlisting}[caption={QEMU 译时与运行时路径（Python 伪代码）。译时把 guest 语义转换成宿主代码，运行时复用；依据源码与开发文档整理的机制示意，并非实测 trace。\cite{source3}},label={lst:tcg-loop}]
def cpu_exec_loop(cpu, pc):
    """运行时主循环：翻译只在缓存未命中时发生"""
    while True:
        tb = tb_cache.lookup(pc, cpu.key_state)
        # ↑ 键 = PC + 翻译相关 CPU 状态（特权级、段基址等）
        if tb is None:                            # 未命中 → 进入译时路径
            tb = tb_gen_code(pc, cpu)             #   现场翻译一个 TB
            tb_cache.insert(pc, cpu.key_state, tb)  # 写入缓存
        pc = tb.run(cpu)                          # 命中 → 直接执行宿主代码
        # 已链接的 TB 在 run() 内部直接跳转，不回到本循环；
        # 链尾或间接跳转（jalr / 返回）才查表回到这里


def tb_gen_code(pc, cpu):
    """译时路径：一个 TB = 译到下一个控制转移为止"""
    ir = []
    while True:
        insn = fetch_decode(pc)               # 取指 + decodetree 模式匹配
        if is_control_transfer(insn):         # jal / beq / ecall：
            ir.append(emit_exit(insn))        #   出口占位、封块
            break
        ir.extend(trans_func[insn.opcode](insn))
        # CSR/MMIO 等敏感指令在 trans_func 内变成一次 helper 调用：
        # 往 IR 流追加一条 call op，语义留在 C 层；
        # helper 调用本身不规定 TB 边界（见 1.3 节）
        pc += insn.length
    return backend_emit(ir)   # 优化 → 寄存器分配 → 发射宿主机器码
\end{lstlisting}

代码块~\ref{lst:tcg-loop} 把上述过程分成译时与运行时两侧：译时是"RISC-V 解码→TCG IR→优化与寄存器分配→x86-64 发射"的流水线，产物放入 TB 缓存；运行时则以 PC 与 CPU 状态查缓存，命中直接执行、未命中回译时，两条路径由块链接与调度循环衔接。需要强调的是，示意中的宿主代码代表语义等价实现，而不是固定的逐指令替换模板。

guest CPU 的机器状态首先表现为 QEMU 进程内的一组 C 结构。RISC-V 的通用寄存器、PC、CSR 与特权级都保存在 \code{CPURISCVState} 中，guest 的当前特权级只是该结构里的一个字段，它与承载 QEMU 进程的宿主线程实际运行在哪个特权级毫无对应关系：guest 的 M 态操作由生成代码或 helper 实现，而不会因此把 QEMU 线程提升到宿主内核态。QEMU 需要宿主服务时仍可通过系统调用进入宿主内核，这与 guest 的特权级变化无关。这是软件虚拟化与 P\&G 直接执行路径最根本的分野，guest 指令不会被宿主硬件执行，因此跨 ISA 的 QEMU 模拟并不满足 P\&G 关于统计上占主导的 guest 指令由真实处理器直接执行的效率要求\cite{source1,source3}。其状态载体与特权级表示如代码块~\ref{lst:cpuriscvstate} 所示：

\begin{lstlisting}[caption={RISC-V guest CPU 状态的主要字段（QEMU target/riscv/cpu.h，节选）},label={lst:cpuriscvstate}]
typedef struct CPUArchState {          /* = CPURISCVState */
    uint64_t gpr[32];                  /* 通用寄存器 */
    uint64_t fpr[32];                  /* 浮点寄存器 */
    uint64_t pc;                       /* 程序计数器 */
    privilege_mode_t priv;             /* 当前特权级：PRV_U/PRV_S/PRV_M */
    uint64_t mstatus;                  /* 特权级状态：MIE、MPP 等 */
    uint64_t mie, mip;                 /* 中断使能 / pending */
    uint64_t satp;                     /* 页表基址寄存器 */
    uint64_t stvec;                    /* S 态异常入口 */
    bool virt_enabled;                 /* H 扩展虚拟化状态 */
    ...
} CPURISCVState;
\end{lstlisting}

一条 helper 的典型执行流程可用 \code{csrrw} 走通，从声明、翻译时发出调用、到运行时代码执行，如代码块~\ref{lst:csrrw-flow} 所示：

\begin{lstlisting}[caption={csrrw helper 的三段流程：声明、翻译时、运行时（节选）},label={lst:csrrw-flow}]
/* 1. 声明：target/riscv/helper.h —— helper 的"函数签名" */
DEF_HELPER_4(csrrw, tl, env, int, tl, tl)

/* 2. 翻译时：target/riscv/tcg/insn_trans/trans_rvi.c.inc */
static bool do_csrrw(DisasContext *ctx, int rd, int rc,
                     TCGv src, TCGv mask)
{
    TCGv dest = dest_gpr(ctx, rd);
    TCGv_i32 csr = tcg_constant_i32(rc);

    translator_io_start(&ctx->base);               /* 标记副作用开始 */
    gen_helper_csrrw(dest, tcg_env, csr, src, mask); /* 追加 call op */
    gen_set_gpr(ctx, rd, dest);   /* helper 返回的旧 CSR 值写回 rd */

    return do_csr_post(ctx);      /* 封块收尾 */
}

/* 3. 运行时：target/riscv/tcg/op_helper.c —— 按宿主 ABI 编译的普通 C 调用 */
target_ulong HELPER(csrrw)(CPURISCVState *env, int csr,
                           target_ulong src, target_ulong write_mask)
{
    target_ulong val = 0;
    RISCVException ret = riscv_csrrw(env, csr, &val, src, write_mask, GETPC());
    if (ret != RISCV_EXCP_NONE)
        riscv_raise_exception(env, ret, GETPC());     /* 转入 guest 异常路径 */
    return val;
}
\end{lstlisting}

翻译时函数体的两处注释在正文展开。\code{translator_io_start}（\path{accel/tcg/translator.c}）标记"有副作用的地方开始了"：把 \code{db->is_jmp} 从 \code{DISAS_NEXT} 改成 \code{DISAS_TOO_MANY}，保证这条指令是当前 TB 的最后一条——副作用（helper 可能抛异常）会推翻跨点的翻译假设，切断最安全，也保证后续 \code{gen_update_pc} 把 PC 写到正确的值。\code{gen_helper_csrrw} 是核心行：往 IR 流"追加"一条 call op——注意是追加，不是现在调用；后端编译这段 IR 时，它才变成按宿主 ABI 摆好参数、call 进 \code{helper_csrrw} 的机器码；\code{tcg_env} 就是 env 指针（guest CPU 状态）。

\code{do_csr_post} 的收尾做三件事：

\begin{enumerate}[label=\arabic*.,leftmargin=2.5em,itemsep=2pt]
\item \code{decode_save_opc} 把本指令的位模式记进 \code{insn_start} 参数（异常回卷还原 tval 用）；
\item \code{gen_update_pc} 把 \code{env->pc} 更新到下一条指令（helper 抛异常时 mepc 要指向正确位置）；
\item \code{exit_tb} 发出退出并设 \code{DISAS_NORETURN} 封块。
\end{enumerate}

为什么写一个 CSR 就要封块？因为 CSR 写可以改 \code{mstatus.MIE}（开关中断）、改特权级、改 \code{satp}（换页表）——刚刚为后续指令做出的翻译假设都可能作废，切断最安全。\code{riscv_csrrw()}（\path{target/riscv/tcg/csr.c}）先用 \code{riscv_csrrw_check()} 检查 CSR 是否实现、最低特权级与只读属性，再经 \code{riscv_csrrw_do64()} 按 CSR 表完成读改写；任一步失败都会返回 \code{RISCV_EXCP_*}，helper 随即转入 guest 异常路径。整个过程中，guest 的“M 态操作”只是对 \code{env->priv} 字段值（\code{PRV_M = 3}）的比较与更新，没有任何宿主特权级参与。

回到译时路径的细节。QEMU 把一条 RISC-V 指令先解码，再表达为与目标无关的中间表示 TCG IR，最后才由宿主后端发射 x86-64 代码。这条三级路径是 QEMU 可移植性的来源：前端只负责把某种 guest ISA 的语义翻译成少量、正交的 TCG 操作，宿主后端只负责把 TCG 操作翻译成某一种宿主指令，两者之间用 IR 解耦，新增一个 guest 或 host 都无需改写另一端。TCG 操作包括临时变量的移动、整数算术（如 \code{add}、\code{and}）、分支、访存以及 \code{call} 等\footnote{QEMU: \path{include/tcg/tcg-opc.h}。}。其中 \code{add} 的译码项落入 \path{target/riscv/insn32.decode} 的 R 型格式，转译函数 \code{trans_add} 最终调用 \code{tcg_gen_add_tl} 发出一次目标宽度加法\footnote{QEMU: \path{target/riscv/tcg/insn_trans/trans_rvi.c.inc}，\code{trans_add()}。}。

翻译以基本块为单位逐块进行：翻译器从当前 PC 出发译到控制转移为止，封出一个 TB。代码块~\ref{lst:rv-src} 是一段最小 RISC-V 程序——经 riscv64-linux-gnu 工具链汇编、链接为静态 ELF，运行退出码 49，与其语义一致（t0=3、t1=4、内存变量 var=42，累加为 49）。它是一条从 \code{_start} 到 \code{ecall} 的直线指令流：三条算术、一条 CSR 读、一次访存、最后退出。

\begin{lstlisting}[caption={guest 程序源码（RISC-V 汇编，riscv64-linux-gnu as+ld 产出静态 ELF，qemu-riscv64 运行退出码 49）},label={lst:rv-src}]
    .section .text
    .globl _start
_start:
    addi t0, zero, 3        # 纯算术：t0 = 3
    addi t1, zero, 4        # 纯算术：t1 = 4
    add  t2, t0, t1         # 纯算术：t2 = t0 + t1
    rdcycle t5              # CSR 读：语义压不进 IR，退 helper
    la   t3, var            # 地址计算（auipc + addi）
    lw   t4, 0(t3)          # 访存：t4 = *var
    add  a0, t2, t4         # 算术 + 访存结果
    li   a7, 93             # 退出系统调用号
    ecall                   # 退出
    .section .data
var:
    .word 42
\end{lstlisting}

代码块~\ref{lst:rv-rdcycle} 与代码块~\ref{lst:rv-ecall} 是该程序在 \code{qemu-riscv64 -d in\_asm,op} 下的真实翻译日志节选（每条 guest 指令前有一行 \code{----} 分隔）。注意一件事：guest 侧本来是一条没有分支的直线流，却被翻译成两个 TB——第 4 条 \code{rdcycle} 是 CSR 读，它的语义压不进 IR 的百余种 op，前端退而生成一次对 helper 的调用，并在调用点之后立即封块；\code{rdcycle} 之后的指令落在下一个 TB。这不是因为控制流变了，而是因为 CSR 访问可能产生副作用，跨点的翻译假设不再可靠。

代码块~\ref{lst:rv-rdcycle} 是第一个 TB 的翻译。三条算术近乎 1:1：\code{addi t0,zero,3} 只产生一条 \code{add_i64}，guest 指令与 IR 一一对应；操作数 \code{x5/t0} 是虚拟寄存器 temp，不绑定任何宿主寄存器——guest 寄存器的“家”在 env 结构里，装入与写回由后端按块内活跃性决定。\code{rdcycle} 则变成 \code{call csrr}（CSR 号 \code{0xc00} 作为立即数传入），随后 \code{mov_i64 pc} 更新 PC、\code{exit_tb} 封块——helper 调用点就是 TB 边界。

\begin{lstlisting}[caption={TB 1：rdcycle 退化为 helper 调用并在调用点封块（qemu-riscv64 -d in\_asm,op 输出节选）},label={lst:rv-rdcycle}]
IN: 0x000100e8:  428d      addi     t0,zero,3
    0x000100ea:  4311      addi     t1,zero,4
    0x000100ec:  006283b3  add      t2,t0,t1
    0x000100f0:  c0002f73  rdcycle  t5

OP:  （TB 1）
 add_i64 x5/t0,$0x0,$0x3                <- addi：纯算术，1:1
 add_i64 x6/t1,$0x0,$0x4
 add_i64 x7/t2,x5/t0,x6/t1
 call csrr,$0x0,$1,x30/t5,env,$0xc00    <- rdcycle：CSR 读退 helper
 mov_i64 pc,$0x100f4
 exit_tb $0x0                           <- helper 调用后封块
\end{lstlisting}

第二个 TB 从 \code{rdcycle} 之后的地址计算继续：\code{auipc} 的翻译只有一条 \code{mov_i64}（目标地址在译时已确定），\code{lw} 则展开成地址计算加一条 \code{qemu_ld_i64}——后端把它展开成软件 TLB 检查序列，快路径直接访存、慢路径转 MMU helper，这是软件虚拟化开销的大头。程序结尾的 \code{ecall}（代码块~\ref{lst:rv-ecall}）是另一种 helper 形态：

\begin{lstlisting}[caption={TB 2 末尾：ecall 退化为 noreturn helper（qemu-riscv64 -d in\_asm,op 输出节选）},label={lst:rv-ecall}]
IN: 0x00010104:  05d00893  addi     a7,zero,93
    0x00010108:  00000073  ecall

OP:  （TB 2 末尾）
 add_i64 x17/a7,$0x0,$0x5d
 mov_i64 pc,$0x10108
 call raise_exception,$0x8,$0,env,$0x8  <- ecall：noreturn helper
\end{lstlisting}

对比两种 helper 形态：\code{rdcycle} 的 \code{call csrr} 会返回（读出的值经输出参数写回 \code{x30/t5}），封块是转译函数出于安全的自选动作；\code{ecall} 的 \code{call raise_exception} 声明为 noreturn，它不返回，guest 的“陷入”被压缩成一次不返回的 C 函数调用。两种形态合起来说明：语义压不进 IR 的操作一律退 helper，但 helper 本身不是 TB 边界——封不封块由转译函数按语义决定。IR 里没有任何 RISC-V 特性：guest 语义被压平成一批通用 op，这就是 m 个前端与 n 个后端只需 m+n 个模块的原因。

翻译产物组织成 Translation Block（TB），即在一组 CPU 状态假设下生成的可独立查找、链接和失效的代码单元。TB 不等同于编译器定义中的严格单基本块，内部可以出现标号与控制流。主执行循环 \code{cpu_exec_loop} 按当前 PC 与翻译相关 CPU 状态查找 TB，未命中才调用 \code{tb_gen_code} 现场翻译，命中则直接执行已缓存的宿主代码\footnote{QEMU: \path{accel/tcg/cpu-exec.c}，\code{cpu_exec_loop()}。}。TB 的有效性不只用 PC 作键：TB 记录了自己成立时假设不变的 CPU 状态（例如特权级、若干段基址），状态变化后必须重新查找匹配当前假设的 TB，必要时才生成新的代码。发射之前，TCG 还要对整块 IR 做优化、可达性分析与活跃性分析，再为临时变量分配宿主寄存器，最后由后端生成机器码——这些步骤在翻译时付出；运行时仍要执行转换后的指令、状态访问与 helper，缓存主要摊薄重复翻译的成本\cite{source3}。为了减少每次回到调度循环的往返，QEMU 提供直接块链接：\code{goto_tb}/\code{exit_tb} 机制在首次跳转后把跳转槽修补成目标 TB 的宿主地址，之后直接跳转；\code{lookup_and_goto_ptr} 则用 helper 查表、再 \code{goto_ptr} 分支\cite{source3}。其中，\code{goto_tb} 直接链接要求目标在翻译时可确定，并满足同页等映射安全约束；运行时查找目标的间接跳转不应与这种直接链接混为一谈。

helper 是表达 guest 复杂语义的手段，而不是 TB 的边界标记：一个 TB 内部可以多次调用 helper，helper 返回后继续执行后续宿主代码。guest 异常同样存在，只是由软件实现。以本章 RISC-V 全系统模式下的访存异常为例：guest 执行 \code{lw t0, 0x4000(t1)} 访一个未映射或权限不足的地址，整条路径如代码块~\ref{lst:mmu-fault} 所示。

\begin{lstlisting}[caption={lw 触发 guest 页错误的处理链路（QEMU target/riscv/tcg/cpu\_helper.c，节选并注释）},label={lst:mmu-fault}]
/* (1) 翻译代码里的 MMU helper 判定翻译或权限失败 */
static void raise_mmu_exception(CPURISCVState *env, target_ulong address,
                                MMUAccessType access_type, ...)
{
    switch (access_type) {
    case MMU_DATA_LOAD:
        /* 按访问类型设置异常索引：数据 load 是页错误 */
        cs->exception_index = RISCV_EXCP_LOAD_PAGE_FAULT;
        break;
    ...
    }
    env->badaddr = address;   /* 出错地址暂存，供异常处理时装入 stval */
}

/* 调用点：失败后恢复该指令的状态、退出生成代码，回到公共执行循环 */
raise_mmu_exception(env, address, access_type, ...);
cpu_loop_exit_restore(cs, retaddr);

/* (2) 公共执行循环兑现这条 guest 异常：riscv_cpu_do_interrupt */
void riscv_cpu_do_interrupt(CPUState *cs)
{
    ...
    if (mode == PRV_S) {                 /* 在 S 态处理 */
        env->scause = cause;             /* 异常号写 scause */
        env->sepc   = env->pc;           /* 当前 PC 存 sepc */
        env->stval  = tval;              /* tval 来自 env->badaddr */
        env->pc = (env->stvec >> 2 << 2);/* 跳 stvec 指向的 guest 处理入口 */
        riscv_cpu_set_mode(env, PRV_S, virt);
    }
    ...
}
\end{lstlisting}

于是 guest 内核看到的是它自己的页错误机制——\code{scause}/\code{stval}/\code{sepc}/\code{stvec} 一个不缺——与真实 RISC-V 处理器上发生的事一模一样。这条路径的“退出”不是 guest 可见的陷阱，而是 QEMU 内部生成代码到解释循环的一次控制转移，随后 guest 层面的异常才被模拟出来。这里不能把其他 ISA 的除法异常，或 user-mode 的宿主信号处理，直接套入本章路径\cite{source3}。

\subsection{内存虚拟化：软件MMU、软件TLB与物理地址空间}\label{sec:software-mem}

guest 的每一次访存都必须经过一段软件地址翻译，同时保留 guest 自己页表与权限的语义。地址关系是三层：guest 虚拟地址经 guest 页表得到 guest 物理地址，再落到 QEMU 进程中的宿主虚拟地址；最后由宿主操作系统与宿主 MMU 完成宿主虚拟地址到宿主物理地址的映射——这最后一段不是 QEMU 模拟的 guest 页表阶段，QEMU 并不干预它。

guest 操作系统负责建立它自己的页表，QEMU 的 RISC-V MMU 实现则在每次地址翻译时读取这些页表、逐级检查权限与访问类型，得到 guest 物理地址，或在违反规则时生成一条 guest 可见的访存异常；对 RISC-V 而言，访问类型还须与 PMP 规则再做一次物理内存保护检查，二者都通过才得到最终可用的 guest 物理地址\footnote{QEMU: \path{target/riscv/tcg/cpu_helper.c}，\code{get_physical_address()}。}。为把这条慢路径从每个访存的热循环里剥离，QEMU 用软件 TLB 缓存翻译结果：TCG 生成的访存代码先检查缓存，RAM 命中时直接以缓存的宿主地址偏移完成读写，未命中才进入 MMU helper 补表\footnote{QEMU: \path{accel/tcg/cputlb.c}，\code{tlb_set_page_full()}；\path{target/riscv/tcg/cpu_helper.c} 中的地址翻译。}。这里必须区分两种"未命中"：软件 TLB miss 通常只是缓存里没有这一页，helper 重走 guest 页表后补上即可继续；而 guest 页表本身不合法时产生的是 guest page fault，会被 \code{raise_mmu_exception} 转成 guest 可见异常。前者是 QEMU 内部的性能事件，后者是模拟出来的架构事件，二者不能混同。

翻译结果之上是物理地址空间的建模。QEMU 用无环的 \code{MemoryRegion} 图与 \code{AddressSpace} 对象描述整台机器的地址布局：RAM 区域有宿主内存后备，读写直达宿主缓冲区；MMIO 区域没有宿主内存，每次访问都会转交对应设备模型处理；此外还有 ROM、ROM 设备、别名与容器等类型，用于表达只读固件、可动态改写的窗口等语义\footnote{QEMU: \path{docs/devel/memory.rst}。}。RAM 与 MMIO 在 TLB 填充阶段就被区分：前者缓存了"guest 地址到宿主内存的偏移"，后者设置 MMIO 等慢路径标志，并不提供可直接读写的 RAM 后备地址，使访存转向设备回调\cite{source3}。这种分流既决定了性能（RAM 快、MMIO 慢），也把内存虚拟化自然接到设备虚拟化上。

两类缓存解决的是两类不同的失效问题：地址翻译变化要求维护软件 TLB，代码内容变化要求失效对应的 TB。软件 TLB 根据 guest 虚拟页及 MMU 上下文选择表项；TB 的物理代码页关联则用于识别被修改的代码并维护失效关系，执行入口另有 PC 相关的快速查找。不能把这些结构统称为物理地址索引的缓存。系统模拟下，软件 MMU 的写入追踪及代码页关联用于失效受影响的 TB，并撤销相应直接链接\footnote{QEMU: \path{accel/tcg/cputlb.c}，\code{tlb_index()}；\path{docs/devel/tcg.rst}，代码失效说明。}\cite{source3}。

\begin{figure}[tbp]
\centering
\begin{tikzpicture}[x=1cm,y=1cm]
\node[box,text width=2.3cm] (va) at (0,0) {guest VA\\取指 / load / store};
\node[box,text width=2.3cm] (tlb) at (3.7,0) {软件 TLB\\检查标签与权限};
\node[greenbox,text width=2.6cm] (fast) at (8.0,0) {RAM 快路径\\由缓存偏移得到 HVA};
\draw[wire] (va)--(tlb); \draw[wire] (tlb)--node[above,font=\footnotesize]{RAM 命中}(fast);
\node[amberbox,text width=2.3cm] (mmu) at (3.7,-1.8) {guest MMU 模型\\查页表与权限};
\node[amberbox,text width=2.3cm] (fault) at (0,-1.8) {guest 访存异常\\交付 guest handler};
\node[box,text width=2.9cm] (as) at (8.0,-1.8) {guest PA\\AddressSpace / MemoryRegion};
\draw[wire] (tlb)--node[left,font=\footnotesize]{miss}(mmu);
\draw[wire] (mmu)--node[above,font=\footnotesize]{失败}(fault);
\draw[wire] (mmu)--node[above,font=\footnotesize]{成功}(as);
\node[greenbox,text width=2.6cm] (ram) at (6.0,-3.7) {RAM 后备内存\\HVA + 页内偏移};
\node[amberbox,text width=2.6cm] (io) at (10.0,-3.7) {MMIO\\调用设备读写回调};
\draw[wire] (as.south)--++(0,-0.45)-|(ram.north);
\draw[wire] (as.south)--++(0,-0.45)-|(io.north);
\node[note] at (4.8,-4.8) {有效映射可填入软件 TLB；宿主 MMU 另行管理 HVA 到 host PA 的转换。};
\end{tikzpicture}
\caption[软件 MMU 的访存路径]{软件 MMU 的访存路径。软件 TLB miss 触发补充转换，不必然产生 guest 页故障；MMIO 和 RAM 共享 guest 物理地址空间，但执行路径不同。\cite{source3}}
\label{fig:softmem}
\end{figure}


图~\ref{fig:softmem} 给出一次 load 的分支全景：RAM 快路径命中时使用缓存地址，miss 时进入 guest MMU 检查映射；合法 RAM 映射可补缓存，MMIO 仍走设备回调，翻译或权限失败则按架构类型生成 guest fault；翻译得到的 guest 物理地址再按 \code{MemoryRegion} 分流为 RAM 直通或 MMIO 回调。整条路径说明：软件虚拟化的访存开销集中在 TLB 命中检查与偶发的翻译慢路径上，而语义正确性依赖对页表与权限的忠实解释。

\subsection{中断虚拟化：从设备事件到guest异常入口}\label{sec:software-irq}

没有任何真实 RISC-V 中断控制器的条件下，异步中断靠一条贯穿软件各层的事件链实现，它从设备状态一路走到 guest 的中断处理程序。事件源可以是串口收到字符、块设备完成请求，也可以是虚拟定时器到期，它们都表现为设备模型内部的一次状态更新。控制器层由 PLIC 与 CLINT 两个模型承担：PLIC 维护外部中断源的优先级、使能与 pending 位，并提供 claim/complete 接口；CLINT 提供所选平台的本地定时器与软件中断能力\cite{source3}。

通知路径是分层的。设备模型通过 \code{qemu_irq} 电平把事件递给控制器，PLIC 在 \code{sifive_plic_update} 里按优先级、使能、未 claim 三个条件逐 hart 挑选应上报的中断，再用 \code{qemu_set_irq} 拉高该 hart 的外部中断输入\footnote{QEMU: \path{hw/intc/sifive_plic.c}，\code{sifive_plic_update()}。}。中断输入的变化最终反映到 CPU 的 pending 状态：\code{riscv_cpu_update_mip} 更新 \code{mip} 寄存器并调用 \code{cpu_interrupt} 置起 \code{CPU_INTERRUPT_HARD}\footnote{QEMU: \path{target/riscv/tcg/cpu_helper.c}，\code{riscv_cpu_update_mip()}。}。此时若 vCPU 正忙于执行翻译代码，还需要 \code{tcg_handle_interrupt} 在非本线程上下文中调用 \code{qemu_cpu_kick} 唤醒或打断该线程\footnote{QEMU: \path{accel/tcg/tcg-accel-ops.c}，\code{tcg_handle_interrupt()}。}。

投递路径则是架构层面的决策。返回公共执行循环后，\code{cpu_handle_interrupt()} 按需调用 RISC-V 目标的 \code{riscv_cpu_exec_interrupt}，它读取 pending、屏蔽位、委托关系与当前特权级，判断该 hart 此刻是否应接收某个中断；应接收时置起异常索引并调用 \code{riscv_cpu_do_interrupt}，由后者记录架构规定的异常返回信息、更新异常相关 CSR、切换特权级并改写 PC，进入 guest 的中断处理入口\footnote{QEMU: \path{target/riscv/tcg/cpu_helper.c}，\code{riscv_cpu_exec_interrupt()}。}。把事件置 pending、打断翻译代码执行、进入架构异常这三步分开是理解本节的关键：置 pending 是设备与控制器的状态变化，kick 是宿主线程调度层面的打断，架构投递才是 guest 状态机的一次转移。TB 链仍须响应异步退出请求。通常，\code{gen_tb_start()} 在生成代码入口插入一个轻量退出检查，通知路径改变相应状态，使后续代码返回公共循环，再做完整的架构中断判断；\code{CF_NOIRQ} 等特殊执行设置可以抑制该检查。因而“检查是否退出”与“完整仲裁并注入 guest 中断”是两层工作，不能用每 TB 全量轮询或完全不检查来概括\footnote{QEMU: \path{accel/tcg/translator.c}，\code{gen_tb_start()}；\path{accel/tcg/tcg-accel-ops.c}，\code{tcg_handle_interrupt()}。}\cite{source3}。

本地定时器使用 QEMU 的虚拟时钟作为软件时间源。所选 CLINT 兼容布局复用 ACLINT 的定时器实现：读取时间时将 \code{QEMU_CLOCK_VIRTUAL} 换算为设备计数，写比较器时更新到期事件；比较条件已满足则立即置起定时器中断，否则在回调中置位\footnote{QEMU: \path{hw/intc/riscv_aclint.c}，\code{cpu_riscv_read_rtc_raw()}、\code{riscv_aclint_mtimer_write_timecmp()} 与 \code{riscv_aclint_mtimer_cb()}。}。这种计时模型不等于默认保证确定性重放。宿主唤醒机制与 guest 中断身份仍是两个不同层面的对象\cite{source3}。

\begin{figure}[tbp]
\centering
\begin{tikzpicture}[x=1cm,y=1cm]
\node[box,text width=2.3cm] (dev) at (0,0) {串口模型\\接收 / 状态更新};
\node[box,text width=2.3cm] (plic) at (3.3,0) {PLIC 模型\\pending / 优先级};
\node[box,text width=2.3cm] (pending) at (6.6,0) {hart 中断输入\\CPU pending 状态};
\node[amberbox,text width=2.3cm] (kick) at (9.9,0) {通知执行线程\\退出请求 / kick};
\draw[wire] (dev)--(plic); \draw[wire] (plic)--(pending); \draw[wire] (pending)--(kick);
\node[box,text width=2.9cm] (timer) at (6.6,1.65) {CLINT 定时器 / 软件中断};
\draw[wire] (timer)--(pending);
\node[amberbox,text width=2.3cm] (check) at (9.9,-2.0) {按架构检查\\屏蔽 / 委托 / 特权};
\node[greenbox,text width=2.3cm] (state) at (6.6,-2.0) {接收中断\\更新 CSR 与 PC};
\node[greenbox,text width=2.3cm] (isr) at (3.3,-2.0) {执行 guest\\中断处理程序};
\node[box,text width=2.3cm] (ack) at (0,-2.0) {PLIC claim / complete\\及设备应答};
\draw[wire] (kick)--(check); \draw[wire] (check)--node[above,font=\footnotesize]{可接收}(state);
\draw[wire] (state)--(isr); \draw[wire] (isr)--(ack);
\draw[dwire] (ack.north)--node[midway,above,sloped,font=\footnotesize]{更新控制器状态}(plic.south);
\end{tikzpicture}
\caption[QEMU 中断生成与投递]{QEMU 中断生成与投递。外部中断以 UART/PLIC 为例；本地定时器和软件中断不经过 PLIC。置 pending、通知线程和 guest 接收中断是不同动作。\cite{source3}}
\label{fig:softirq}
\end{figure}


图~\ref{fig:softirq} 以串口接收为例给出时序：字符到达→PLIC 置 pending 并选优→hart 中断输入拉高→TCG 退出或 kick→CPU 判断接收→guest handler 执行→claim/complete 回写控制器。这条链把三个层次的介入点都标了出来，说明软件中断虚拟化的代价主要分布在控制器状态维护与每次投递判断上。

\subsection{设备虚拟化：MMIO、软件DMA与后端服务}\label{sec:software-dev}

一个设备除了寄存器本身，还必须模拟三个方向的行为：CPU 到设备的控制路径、设备到内存的数据路径、设备回 CPU 的完成通知。

控制路径由 MMIO 分发实现。guest 对设备寄存器所在地址的读写，先经物理地址空间定位到 MMIO 区域，再转交 \code{MemoryRegionOps} 里注册的回调。以 NS16550 兼容串口为例，\code{serial_mm_write} 把写值裁到 8 位后送入 \code{serial_io_ops}，寄存器写会改变内部状态并可能触发字符后端动作，寄存器读则返回当前状态\footnote{QEMU: \path{hw/char/serial-mm.c}，\code{serial_mm_write()}。}。也就是说，MMIO 分发把一次 guest 设备寄存器写转换为对设备状态机的调用，而不是先对某个真实 MMIO 地址进行宿主写操作，guest 看到的只是一个"存在寄存器"的串口，后端连的是字符设备、管道还是文件由 QEMU 配置决定。

数据路径以 virtio-blk 展开。guest 通过 virtqueue 提交描述符链，QEMU 在 \path{virtio_blk_handle_vq} 里取请求、解析描述符，把块后端处理结果写回 guest 内存，最后在 \path{virtio_blk_req_complete} 里写状态字节、\code{virtqueue_push} 归还描述符并 \code{virtio_notify} 通知 guest\footnote{QEMU: \path{hw/block/virtio-blk.c}，\code{virtio_blk_handle_vq()}、\code{virtio_blk_req_complete()}。}。这里的数据搬运是软件 DMA：设备通过自己的 \code{AddressSpace} 访问 guest 内存，直接使用 guest 物理地址到宿主缓冲区的映射，并不自动套用某个 guest 进程的 CPU 页表；只有当配置了 vIOMMU 时，设备地址才会额外经过一次设备地址转换\cite{source3}。区分"软件 DMA 走设备地址空间"与"CPU 访存走 CPU 页表"是理解设备与内存两层如何分工的要点。完成通知则是回接上一节：请求完成后设备更新队列状态，需要时拉起中断，走 PLIC 与架构投递回到 guest。

最后要区分两种设备接口。模拟现有硬件接口（如 NS16550、PCI 设备）与 virtio 这类半虚拟化设备接口，都能由 QEMU 的软件后端提供服务，差别只在 guest 侧协议与性能；virtio 不是 CPU 硬件虚拟化指令，也不要求 CPU 必须用 KVM 执行，它只是约定了一套更利于软件实现高效 I/O 的接口。具体到 virtio-blk，virtqueue 组织请求与完成信息。以 split virtqueue 为例，guest 写好描述符链、更新 avail 环并按通知规则发出 kick，设备侧读描述符、执行后端、把结果写回 guest 内存、再写 used 环并按协商的通知规则决定是否产生中断，二者通过共享内存环而非真实总线握手来传递请求与完成状态。两者在"由谁实现设备、guest 如何发现与驱动"上不同，在"都是软件模型"这一点上相同。

\subsection{能力与成本}\label{sec:software-cost}

把四个维度并置，才能看出软件方案"能做什么"与"为什么慢"是同一件事的两面。翻译代码让 guest 计算得以进行，软件地址转换让 guest 访存保持页表语义，事件状态机让异步中断成立，设备模型让 I/O 可用，四者组合起来才足以运行一套完整的 guest 操作系统；缺任何一环，其余环都没有意义。这是软件虚拟化的完整能力边界：它不依赖宿主硬件的任何虚拟化特性，因而可移植性极强——增加目标平台主要涉及相应 CPU/设备模型和宿主执行后端，而不必为每个 guest-host 组合重写整套实现\cite{source2,source3}。

成本则须逐项拆解，而不是笼统写"慢若干倍"。翻译与缓存维护的代价集中在冷路径（首次翻译）与失效路径（映射或代码变化），热路径只付出块链接的跳转与状态检查；访存代价来自每个访存前的软件 TLB 检查，以及 miss 时重走 guest 页表的开销；helper 调用打断了宿主代码的寄存器分配局部性，并在 C 代码与生成代码之间切换上下文；中断与设备 I/O 还引入 vCPU 线程与事件线程之间的同步，以及后端真实 I/O 的等待。这些成本落在不同频率与不同路径上：一个纯计算密集的 guest 主要承受翻译与寄存器分配的摊还成本，一个 I/O 密集的 guest 则主要承受 MMIO 回调、中断往返与后端等待，脱离具体工作负载无法给出统一的性能倍数——本综述亦不采信任何未指明配置与基线的性能数字。

这套方案的边界也就此划清：它以软件实现架构语义，并在翻译执行、地址转换与中断投递中承担额外管理工作；功能等价依赖各模型及其组合的正确性，而不能仅由“纯软件”这一实现方式保证。下一章考察的硬件辅助，正是把这些高频路径上的软件介入下放到处理器与系统结构，但设备功能与管理策略仍需软件完成，硬件并不会自动补上软件方案中"模型"那一半。
